\documentclass[11pt]{article}

\usepackage[a4paper,margin=2.4cm]{geometry}
\usepackage{amsmath,amssymb}
\usepackage{booktabs}
\usepackage{siunitx}
\usepackage[hidelinks]{hyperref}
\usepackage{enumitem}

\newcommand{\tirr}{t_{\mathrm{irr}}}
\newcommand{\tdel}{t_{\mathrm{del}}}
\newcommand{\tmeas}{t_{\mathrm{meas}}}

\title{The A$\to$B handoff:\\
time-decay bookkeeping, source sampling, and $\sigma(\text{range})$}
\author{Method note for review}
\date{\today}

\begin{document}
\maketitle

\begin{abstract}
This note documents the handoff (Stage~B0, \texttt{decay\_sampling/}) that turns
the Stage-A \emph{production} of $\beta^{+}$ emitters into the \emph{measured}
annihilation source seen by a detector, and explains the Monte-Carlo definition
of the figure of merit $\sigma(\text{range})$. Three ideas: (i) the detector
does not see production --- it sees the decays that fall inside its acquisition
window, after a delay, scaled to a clinical dose; (ii) the source is delivered
as a \emph{budget} (how many decays of each isotope, plus a seed), not a
materialised event list; (iii) $\sigma(\text{range})$ is the spread of the
recovered range over many Poisson-noisy repeats of the same measurement --- the
error bar on the single real measurement.
\end{abstract}

\section{Role in the chain}

Stage~A produced \texttt{emitters.csv}: one row per produced emitter, carrying
its annihilation point and isotope. This is the production \emph{shape}
$P_j(\mathbf r)$ at the statistics of the (arbitrary) number of protons we chose
to simulate. The detector (Stage~B) needs a list of annihilation \emph{events}
to fire 511\,keV pairs from. We cannot feed \texttt{emitters.csv} directly,
for two reasons:
\begin{itemize}[nosep]
\item \textbf{Wrong absolute number.} Our file holds the emitters of $\sim10^6$
simulated protons; a clinical $1$\,Gy field is $\sim10^{10}$ protons.
\item \textbf{Wrong isotope mix.} The detector sees only the decays inside its
acquisition window, which starts a delay after irradiation; by then short-lived
$^{15}$O has largely decayed, so the \emph{measured} mix $\neq$ the
\emph{produced} mix.
\end{itemize}
The handoff fixes both: scale to a clinical dose (Sec.~\ref{sec:norm}), then
apply radioactive decay and the acquisition timing (Sec.~\ref{sec:decay}).

\section{Time-decay model}
\label{sec:decay}

Each species $j$ has decay constant $\lambda_j=\ln2/T_{1/2,j}$ and is assumed to
decay $100\%$ by $\beta^{+}$. For a delivery of duration $\tirr$ at constant
production rate $K_j=P_j/\tirr$, the activity is
\begin{align}
A_j(t) &= K_j\bigl(1-e^{-\lambda_j t}\bigr), & 0\le t\le \tirr
       &\quad\text{(build-up)},\\
A_j(t) &= K_j\bigl(1-e^{-\lambda_j \tirr}\bigr)\,e^{-\lambda_j (t-\tirr)},
       & t>\tirr &\quad\text{(decay)}.
\end{align}
A scanner acquires over the window $[t_0,t_1]$ with $t_0=\tirr+\tdel$ and
$t_1=t_0+\tmeas$. The number of decays collected is
$N_j=\int_{t_0}^{t_1}A_j(t)\,dt$, which evaluates to the \textbf{three-factor
expression} (continuous / cyclotron limit, the relevant case --- post-beam
acquisition makes the pulse structure irrelevant):
\begin{equation}
\boxed{\;
N_j = P_j\;
\underbrace{\frac{1-e^{-\lambda_j \tirr}}{\lambda_j \tirr}}_{\text{build-up}}\;
\underbrace{e^{-\lambda_j \tdel}}_{\text{transport}}\;
\underbrace{\bigl(1-e^{-\lambda_j \tmeas}\bigr)}_{\text{window}}\;}
\label{eq:nmeas}
\end{equation}
The factors are: how many nuclei survive to the end of beam (production and
decay during $\tirr$); decay during the transport delay $\tdel$; the fraction of
survivors that decay during the acquisition $\tmeas$.

\paragraph{Operating point (Parodi standard).} $\tirr=\SI{60}{s}$,
$\tdel=\SI{300}{s}$, $\tmeas=\SI{1800}{s}$ (cyclotron, offline). All are
parameters.

\paragraph{Activity curves --- a required output.} Plotting $A_j(t)$ per isotope
from $t=0$ through the acquisition window, with the window shaded, reproduces the
classic activity-vs-time picture (CRYSP paper Fig.~1 / Parodi): build-up during
$\tirr$, then isotope-specific decay. It shows visually why a short delay is
decisive --- at $\tdel\sim$ a few minutes $^{15}$O ($T_{1/2}=\SI{122}{s}$) is
still abundant, whereas at offline delays only $^{11}$C remains. The handoff
emits this plot.

\section{Absolute normalisation}
\label{sec:norm}

Stage~A also scores the dose in the target box, $D_{\mathrm{sim}}$, for the same
$N_{\mathrm{sim}}$ protons (\texttt{run\_meta.csv}: \texttt{target\_dose\_Gy}).
The production for a clinical dose $D$ is then
\begin{equation}
P_j(D) = \mathrm{count}_j \,\frac{D}{D_{\mathrm{sim}}},
\label{eq:pj}
\end{equation}
i.e.\ $P_j(\SI{1}{Gy})=\mathrm{count}_j/D_{\mathrm{sim}}$ (Sec.~\ref{sec:norm}
closes the loop with the Parodi Table~2 cross-check,
\texttt{analysis\_transport/parodi\_cross\_check.py}). Substituting
Eq.~\eqref{eq:pj} into Eq.~\eqref{eq:nmeas} gives the measured decays $N_j$ for a
$D$-Gy field. For the standard scenario the totals are $N\sim10^{7}$--$10^{8}$.

\section{The detector source: sampling}
\label{sec:source}

Stage~B needs $N_j$ annihilation \emph{points} per isotope. The file holds only
$\sim10^4$ points per isotope (the shape); we obtain $N_j$ of them by
\textbf{sampling with replacement} --- drawing points at random (repeats
allowed). Drawing $N_j$ times reproduces the shape $P_j(\mathbf r)$ at the
required count.

\paragraph{Budget, not materialised events.} Writing all $\sim10^{8}$ sampled
points to disk would be GB-scale per realisation. Instead the handoff writes a
\textbf{budget}: the per-isotope counts $N_j$ (and the Poisson draws below) plus
a random seed and the operating point. Stage~B does the sampling \emph{on the
fly} from \texttt{emitters.csv}. With the same file, budget and seed the drawn
set is identical and reproducible, so \textbf{every detector consumes the
identical source} --- guaranteed by determinism, not by sharing a large file.

\section{The figure of merit: $\sigma(\text{range})$}
\label{sec:sigma}

\paragraph{What it is.} The precision with which the measurement recovers the
proton range --- one number in mm. It is the \emph{error bar} on the single real
measurement: a clinic scans once and gets a range, but that value is one draw
from a distribution whose width is $\sigma(\text{range})$. You cannot read the
error bar off one number, so we characterise it by simulation.

\paragraph{Why Poisson.} The measurement counts independent rare events
(detected decays), so the realised integer count of species $j$ fluctuates as
Poisson about the expectation $N_j$ (variance $N_j$). The expectation $N_j$
(Eq.~\eqref{eq:nmeas}) is deterministic; any one acquisition is a Poisson draw
$M_j\sim\mathrm{Poisson}(N_j)$.

\paragraph{The $Z$ realisations.} We repeat the \emph{whole} measurement $Z$
($\sim100$) times. Each realisation $z$:
\begin{enumerate}[nosep]
\item draws $M_j^{(z)}\sim\mathrm{Poisson}(N_j)$ for \emph{every} isotope, with
its own seed;
\item forms the source (those $M_j^{(z)}$ sampled points), runs the detector and
reconstruction;
\item fits the distal fall-off $\to$ a range estimate $R_z$.
\end{enumerate}
After $Z$ realisations, $\{R_1,\dots,R_Z\}$ are independent estimates from
independent noisy draws; their
\begin{equation}
\sigma(\text{range}) = \operatorname{std}\{R_z\}
\end{equation}
is the precision. (Analogy: a scale's precision is characterised once by
repeated tests; thereafter every single weighing carries that $\pm$.) The real
treatment happens once, but between hypothetical repeats only the counting noise
changes --- which we can simulate --- so the spread over $Z$ draws is the
uncertainty on the one real outcome. The detector with the smallest
$\sigma(\text{range})$ at a fixed budget wins.

\paragraph{Tractability (Stage B/C detail).} We do \emph{not} run the Geant4
detector $Z$ times ($Z\times10^{8}$ would be $\sim10^{10}$ pairs). The detector
is simulated \textbf{once} on the full source $\to$ a master coincidence list;
each realisation is a Poisson-sized \emph{resample} of that list, and only the
cheap reconstruction runs $Z$ times. The exact resampling scheme is fixed when
Stage~B is built.

\section{Implementation (\texttt{decay\_sampling/})}

\begin{itemize}[nosep]
\item Inputs: \texttt{emitters.csv}, \texttt{run\_meta.csv}; parameters
$D,\tirr,\tdel,\tmeas,Z$, seed.
\item $\lambda_j$ from \texttt{common/isotopes.py}; $N_j$ from
Eqs.~\eqref{eq:pj},\eqref{eq:nmeas}; $M_j^{(z)}\sim\mathrm{Poisson}(N_j)$.
\item Output: \texttt{sampling\_budget\_<scenario>.csv} --- per isotope
$N_j$ and the $Z$ Poisson draws, plus the operating point and seeds (the budget,
Sec.~\ref{sec:source}).
\item A script emitting the activity-vs-time plot $A_j(t)$ (Sec.~\ref{sec:decay}).
\end{itemize}

\end{document}
